\documentclass[a4paper, 11pt]{article}
\usepackage[utf8]{inputenc}
\usepackage[T1]{fontenc}
\usepackage[italian]{babel}
\usepackage{geometry}
\usepackage{graphicx}
\usepackage{tabularx}
\usepackage{fancyhdr}
\usepackage{hyperref}

% Margini
\geometry{top=3cm, bottom=3cm, left=2.5cm, right=2.5cm}

% Intestazione e piè di pagina
\pagestyle{fancy}
\fancyhf{}
\fancyhead[L]{\small \textit{0x06tudio} -- Verbale di Riunione}
\fancyhead[R]{\small Anno Accademico 2026/2027}
\fancyfoot[C]{\thepage}
\renewcommand{\headrulewidth}{0.4pt}

\begin{document}
\thispagestyle{plain}

\newcommand{\issue}[1]{\href{https://github.com/0x06tudio/docs/issues/#1}{\textit{#1}}}

\begin{minipage}{0.6\textwidth}
    \vspace{-2cm}
    {\huge \textbf{0x06tudio}} \\[0.5em]
    {\large Progetto didattico di Ingegneria del Software} \\[0.3em]
    {\normalsize Università degli Studi di Padova}
\end{minipage}
\begin{minipage}{0.4\textwidth}
    \raggedleft
    \vspace{0pt}
    \includegraphics[width=6cm]{../../media/logo.png}
    \vspace{1cm}
    \small
    \href{mailto:0x06tudio@gmail.com}{0x06tudio@gmail.com}
\end{minipage}

\vspace{0.5cm}

\begin{center}
    {\LARGE \textbf{Verbale Riunione}}
\end{center}

\vspace{0.5cm}


\begin{tabularx}{\textwidth}{l X}
    \textbf{Versione:} & 0.2.0 \\
    \textbf{Data:} & 2026/10/30 \\
    \textbf{Uso:} & Interno \\
    \textbf{Partecipanti (Presenti):} & Nicolò Beltrame, Pietro Brocadello, Alberto Cavalletto, Catalin Chirilov, Davide Colabove, Andrea Maggio \\
    \textbf{Destinatari:} & 0x06tudio \\
    \textbf{Modalità / Luogo:} & Online / Discord \\
    \textbf{Ora inizio:} & 15:30 \\
    \textbf{Ora fine:} & 17:00 \\
    \textbf{Durata totale:} & 1 ora 30 minuti \\
\end{tabularx}
\section*{Ordine del giorno}
\begin{enumerate}
    \item Revisione documenti\vspace{-0.2cm}
    \item Invio lettera di candidatura\vspace{-0.2cm}
    \item Pianificazione per la realizzazione del diario di bordo\vspace{-0.2cm}
    \item Studio delle tecnologie per il capitolato C9 (SpecCheck)\vspace{-0.2cm}
\end{enumerate}
\vspace{-0.5cm}

\section*{Registro delle modifiche}
\begin{tabularx}{\textwidth}{|c|c|X|c|c|}
    \hline
    \textbf{Versione} & \textbf{Data} & \textbf{Descrizione} & \textbf{Verificatore} & \textbf{Autore} \\ \hline
    0.2.0 & 2026/11/03 & Correzioni per nuova candidatura &  & Andrea Maggio \\ \hline
    0.1.0 & 2026/10/30 & Revisione delle documentazioni e invio lettera di candidatura & Nicolò Beltrame & Catalin Chirilov \\ \hline
\end{tabularx}

\newpage
\tableofcontents
\newpage


\section{Resoconto della discussione}

\subsection{Punto 1: Revisione}
Il gruppo ha revisionato i documenti precedentemente redatti, dove sono emersi piccoli errori di contenuto nelle intestazioni e alcune sviste morfosintattiche.  

\subsection{Punto 2: Lettera di candidatura}
Una volta assicurata la correttezza dei documenti, si è proseguito con l'invio della lettera di candidatura ai committenti.

\subsection{Punto 3: Pianificazione diario di bordo}
A seguito di una discussione, il gruppo ha stabilito in dettaglio gli argomenti da presentare per il Diario di Bordo 2 (DdB 2). I punti salienti individuati sono i seguenti:
\begin{itemize}
    \item \textbf{Stato di avanzamento:} incontro con VarGroup, firma dei verbali esterni, stesura e presentazione della candidatura e assegnazione dei ruoli.
    \item \textbf{Obiettivi a breve termine:} implementare le automazioni nella repository e avviare la ricerca per i cruscotti di progetto.
    \item \textbf{Difficoltà riscontrate:} difficoltà relative alla stima dei costi e delle ore necessarie.
    \item \textbf{Dubbi ed incertezze:} individuazione delle risorse da studiare per proseguire i lavori.
\end{itemize}

\subsection{Punto 4: Studio delle tecnologie per il capitolato C9 (SpecCheck)}
In previsione di un'eventuale assegnazione del capitolato prescelto (C9 - SpecCheck), il gruppo ha concordato la necessità di avviare uno studio individuale delle varie tecnologie richieste. A valle di questo studio, verrà redatto un documento apposito che spieghi in modo dettagliato il funzionamento di tali tecnologie e le motivazioni per il loro utilizzo all'interno del progetto.

\newpage


\section{Riepilogo delle decisioni e degli impegni (To-Do)}

\begin{tabularx}{\textwidth}{|c|X|X|}
    \hline
    \textbf{ID Decisione} & \textbf{Descrizione} & \textbf{Motivazione} \\ \hline
    D 4.1 & Approvata revisione finale e invio della lettera di candidatura. & Consenso unanime \\ \hline
    D 4.2 & Definiti i punti chiave da presentare nel Diario di Bordo 2 (DdB 2). & Brainstorming in vista della lezione del 2026/10/02 \\ \hline
    D 4.3 & Approvato lo studio individuale delle tecnologie (Capitolato C9) e la stesura dei relativi documenti esplorativi. & Documenti per facilitare la presa in carico dei nuovi ruoli al momento della rotazione \\ \hline
\end{tabularx}

\vspace{2cm}

\noindent
\begin{tabularx}{\textwidth}{|c|X|c|c|c|}
    \hline
    \textbf{Task ID} & \textbf{Descrizione} & \textbf{Assegnato a} & \textbf{Scadenza} & \textbf{Issue} \\ \hline
    TD 4.1 & Ricerca cruscotto rendicontazione ore & Tutti i membri & 2026/11/03 & \issue{25} \\ \hline
    TD 4.2 & Ricerca cruscotto qualità & Tutti i membri & 2026/11/03 & \issue{26} \\ \hline
    TD 4.3 & Sistemare la repository con automazioni & Tutti i membri & 2026/11/06 & \issue{27} \\ \hline
    TD 4.4 & Studio tecnologie (C9) e stesura documenti esplorativi & Tutti i membri & 2026/11/09 & \issue{28} \\ \hline
\end{tabularx}

\end{document}